% GENERATED FILE. Edit canonical lessons, then run npm run generate:books.
% canonical-source-hash: fnv1a64:1aaee8e1a6f2bf53
% canonical-lessons: HI-C440-dhul, HI-C440-ghirna, HI-C440-thamnaintr, HI-C440-chatnaintr, HI-C440-urana

\chapter{Clouds Gather and Clear}
\label{ch:hi-clouds-gather-and-clear}
\hlchaptermodality{440}

\begin{chapteropening}
\textbf{By the end of this chapter:} \emph{I can say a village was cut off by water, the rain let up, the sky cleared and the wind blew things away.}

Complete the last lesson of chapter 440: I can say a village was cut off by water, the rain let up, the sky cleared and the wind blew things away.
\end{chapteropening}

\section[\texorpdfstring{\hi{धूल} (dhūl)}{धूल (dhūl)}]{\hi{धूल} (dhūl) — dust}

\label{lesson:HI-C440-dhul}



\begin{quote}
Before the new one: say the Hindi for a degree, then the Hindi for to scorch.
\end{quote}

\subsection*{You'll want to know: \hi{धूल}}
\textbf{\hi{धूल}} — \emph{dhūl} — \textquotedblleft{}dust\textquotedblright{}.

\begin{quote}
\hi{आँधी} \hi{के} \hi{बाद} \hi{हर} \hi{चीज़} \hi{पर} \hi{धूल} \hi{थी।}
\end{quote}

\begin{quote}
\emph{ā̃dhī ke bād har cīz par dhūl thī.}
\end{quote}

\begin{quote}
After the dust storm everything was covered in dust.
\end{quote}

Feminine.

\begin{cousinweb}
\textbf{\hi{धूल}} (\emph{dhūl}) comes from Sanskrit \emph{dhūli}, \textquotedblleft{}dust, powder\textquotedblright{}, which CDIAL files under \emph{*dhūḍi} (CDIAL 6835) with Hindi \emph{dhūl}, \textquotedblleft{}dust\textquotedblright{}.
\end{cousinweb}

\begin{grammarlens}[title={What you've built}]
One more word for getting around.
\end{grammarlens}

\subsection*{Your turn}
\begin{itemize}
\raggedright
  \item \emph{Say it:} \emph{dhūl}
  \item \emph{Say it:} \emph{dhūl}, once more
  \item \emph{Say it:} \emph{ā̃dhī ke bād har cīz par dhūl thī.} — After the dust storm everything was covered in dust.
  \item \emph{Recall:} say the Hindi for a degree, then the Hindi for to scorch, then say \emph{dhūl} again
\end{itemize}

\begin{tcolorbox}[breakable,colback=teal!4,colframe=teal!35!black,title={Before you move on}]
What does \hi{धूल} mean? (\textquotedblleft{}Dust\textquotedblright{}.) Say it once more.
\end{tcolorbox}
\section[\texorpdfstring{\hi{घिरना} (ghirnā)}{घिरना (ghirnā)}]{\hi{घिरना} (ghirnā) — to be surrounded; to gather round}

\label{lesson:HI-C440-ghirna}



\begin{quote}
Before the new one: say the Hindi for to scorch, then the Hindi for dust.
\end{quote}

\subsection*{You'll want to know: \hi{घिरना}}
\textbf{\hi{घिरना}} — \emph{ghirnā} — \textquotedblleft{}to be surrounded; to gather round\textquotedblright{}.

\begin{quote}
\hi{बाढ़} \hi{में} \hi{हमारा} \hi{गाँव} \hi{पानी} \hi{से} \hi{घिर} \hi{गया।}
\end{quote}

\begin{quote}
\emph{bāṛh meṁ hamārā gā̃v pānī se ghir gayā.}
\end{quote}

\begin{quote}
In the flood our village was cut off by water.
\end{quote}

Past (it …): \hi{वह} \hi{मैदान} \hi{पेड़ों} \hi{से} \hi{घिरा}, \emph{vah maidān peṛoṁ se ghirā}.

\begin{cousinweb}
\textbf{\hi{घिरना}} (\emph{ghirnā}) goes back to \emph{*ghir}, \textquotedblleft{}go round\textquotedblright{}, a form CDIAL reconstructs and thinks may come from Dravidian (CDIAL 4474); it lists Hindi \emph{ghirnā}, \textquotedblleft{}to be surrounded\textquotedblright{}.
\end{cousinweb}

\begin{grammarlens}[title={What you've built}]
One more word for getting around.
\end{grammarlens}

\subsection*{Your turn}
\begin{itemize}
\raggedright
  \item \emph{Say it:} \emph{ghirnā}
  \item \emph{Say it:} \emph{ghirnā}, once more
  \item \emph{Say it:} \emph{bāṛh meṁ hamārā gā̃v pānī se ghir gayā.} — In the flood our village was cut off by water.
  \item \emph{Recall:} say the Hindi for to scorch, then the Hindi for dust, then say \emph{ghirnā} again
\end{itemize}

\begin{tcolorbox}[breakable,colback=teal!4,colframe=teal!35!black,title={Before you move on}]
What does \hi{घिरना} mean? (\textquotedblleft{}To be surrounded; to gather round\textquotedblright{}.) Say it once more.
\end{tcolorbox}
\section[\texorpdfstring{\hi{थमना} (thamnā)}{थमना (thamnā)}]{\hi{थमना} (thamnā) — to let up; to stop (of rain or wind)}

\label{lesson:HI-C440-thamnaintr}



\begin{quote}
Before the new one: say the Hindi for dust, then the Hindi for to be surrounded.
\end{quote}

\subsection*{You'll want to know: \hi{थमना}}
\textbf{\hi{थमना}} — \emph{thamnā} — \textquotedblleft{}to let up; to stop (of rain or wind)\textquotedblright{}.

\begin{quote}
\hi{दो} \hi{घंटे} \hi{बाद} \hi{बारिश} \hi{थमी।}
\end{quote}

\begin{quote}
\emph{do ghaṇṭe bād bāriś thamī.}
\end{quote}

\begin{quote}
The rain let up after two hours.
\end{quote}

Past (it …): \hi{वह} \hi{आँधी} \hi{थम} \hi{गई}, \emph{vah ā̃dhī tham gaī}.

\begin{cousinweb}
\textbf{\hi{थमना}} (\emph{thamnā}) belongs with Sanskrit \emph{stambhate}, \textquotedblleft{}supports, arrests\textquotedblright{} (CDIAL 13683), which lists Hindi \emph{thāmnā}, \textquotedblleft{}to prop, stop\textquotedblright{}; \hi{थमना} is its intransitive: the rain stops itself.
\end{cousinweb}

\begin{grammarlens}[title={What you've built}]
One more word for getting around.
\end{grammarlens}

\subsection*{Your turn}
\begin{itemize}
\raggedright
  \item \emph{Say it:} \emph{thamnā}
  \item \emph{Say it:} \emph{thamnā}, once more
  \item \emph{Say it:} \emph{do ghaṇṭe bād bāriś thamī.} — The rain let up after two hours.
  \item \emph{Recall:} say the Hindi for dust, then the Hindi for to be surrounded, then say \emph{thamnā} again
\end{itemize}

\begin{tcolorbox}[breakable,colback=teal!4,colframe=teal!35!black,title={Before you move on}]
What does \hi{थमना} mean? (\textquotedblleft{}To let up; to stop (of rain or wind)\textquotedblright{}.) Say it once more.
\end{tcolorbox}
\section[\texorpdfstring{\hi{छँटना} (chãṭnā)}{छँटना (chãṭnā)}]{\hi{छँटना} (chãṭnā) — to clear (of clouds); to thin out}

\label{lesson:HI-C440-chatnaintr}



\begin{quote}
Before the new one: say the Hindi for to be surrounded, then the Hindi for to let up.
\end{quote}

\subsection*{You'll want to know: \hi{छँटना}}
\textbf{\hi{छँटना}} — \emph{chãṭnā} — \textquotedblleft{}to clear (of clouds); to thin out\textquotedblright{}.

\begin{quote}
\hi{बारिश} \hi{के} \hi{बाद} \hi{बादल} \hi{छँट} \hi{गए।}
\end{quote}

\begin{quote}
\emph{bāriś ke bād bādal chãṭ gae.}
\end{quote}

\begin{quote}
After the rain the clouds cleared.
\end{quote}

Past (it …): \hi{वह} \hi{कोहरा} \hi{छँट} \hi{गया}, \emph{vah kohrā chãṭ gayā}.

\begin{cousinweb}
\textbf{\hi{छँटना}} (\emph{chãṭnā}) goes back to \emph{*chaṇṭ}, \textquotedblleft{}scatter\textquotedblright{}, a form CDIAL reconstructs (CDIAL 4970), which lists Hindi \emph{chā̃ṭnā}, \textquotedblleft{}to sift, select\textquotedblright{}; clouds that \hi{छँट} \hi{जाते} are scattered and thinned away.
\end{cousinweb}

\begin{grammarlens}[title={What you've built}]
One more word for getting around.
\end{grammarlens}

\subsection*{Your turn}
\begin{itemize}
\raggedright
  \item \emph{Say it:} \emph{chãṭnā}
  \item \emph{Say it:} \emph{chãṭnā}, once more
  \item \emph{Say it:} \emph{bāriś ke bād bādal chãṭ gae.} — After the rain the clouds cleared.
  \item \emph{Recall:} say the Hindi for to be surrounded, then the Hindi for to let up, then say \emph{chãṭnā} again
\end{itemize}

\begin{tcolorbox}[breakable,colback=teal!4,colframe=teal!35!black,title={Before you move on}]
What does \hi{छँटना} mean? (\textquotedblleft{}To clear (of clouds); to thin out\textquotedblright{}.) Say it once more.
\end{tcolorbox}
\section[\texorpdfstring{\hi{उड़ाना} (uṛānā)}{उड़ाना (uṛānā)}]{\hi{उड़ाना} (uṛānā) — to blow away; to fly (a kite)}

\label{lesson:HI-C440-urana}



\begin{quote}
Before the new one: say the Hindi for to let up, then the Hindi for to clear (of clouds).
\end{quote}

\subsection*{You'll want to know: \hi{उड़ाना}}
\textbf{\hi{उड़ाना}} — \emph{uṛānā} — \textquotedblleft{}to blow away; to fly (a kite)\textquotedblright{}.

\begin{quote}
\hi{हवा} \hi{ने} \hi{मेरी} \hi{टोपी} \hi{उड़ा} \hi{दी।}
\end{quote}

\begin{quote}
\emph{havā ne merī ṭopī uṛā dī.}
\end{quote}

\begin{quote}
The wind blew my cap away.
\end{quote}

Past, with \hi{ने} (he or she …): \hi{उसने} \hi{पतंग} \hi{उड़ाई}, \emph{usne patang uṛāī}.

\begin{cousinweb}
\textbf{\hi{उड़ाना}} (\emph{uṛānā}) is the causative partner of \hi{उड़ना}, \textquotedblleft{}to fly\textquotedblright{}, from Sanskrit \emph{uḍḍīyate}, \textquotedblleft{}flies up\textquotedblright{} (CDIAL 1697), which lists Hindi \emph{uṛānā}, \textquotedblleft{}to cause to fly\textquotedblright{}.
\end{cousinweb}

\begin{grammarlens}[title={What you've built}]
One more word for getting around.
\end{grammarlens}

\subsection*{Your turn}
\begin{itemize}
\raggedright
  \item \emph{Say it:} \emph{uṛānā}
  \item \emph{Say it:} \emph{uṛānā}, once more
  \item \emph{Say it:} \emph{havā ne merī ṭopī uṛā dī.} — The wind blew my cap away.
  \item \emph{Recall:} say the Hindi for to let up, then the Hindi for to clear (of clouds), then say \emph{uṛānā} again
\end{itemize}

\begin{tcolorbox}[breakable,colback=teal!4,colframe=teal!35!black,title={Before you move on}]
What does \hi{उड़ाना} mean? (\textquotedblleft{}To blow away; to fly (a kite)\textquotedblright{}.) Say it once more.
\end{tcolorbox}
